% !TEX root = ../main.tex
% =====================================================================
% CAPITOLO 5 — ORION: ARCHITETTURA
% =====================================================================

\chapter{Orion: architettura}
\label{chap:orion-architettura}

Questo capitolo presenta le scelte architetturali dell'orchestratore Orion, il componente ``conclusivo'' della tesi.
Mentre i Capitoli~\ref{chap:agenti-architettura} e~\ref{chap:agenti-implementazione} hanno descritto la 
suite di agenti verticali come un insieme di servizi indipendenti, 
ciascuno raggiungibile dall'utente attraverso il portale, qui si affronta il problema complementare: 
come consentire all'utente di interagire con \emph{un unico punto di accesso} che instradi automaticamente 
ogni richiesta verso l'agente più competente.

% =====================================================================
\section{Motivazione e visione architetturale}
\label{sec:orion-motivazione}
% =====================================================================

L'approccio ad unico agente orchestratore consente il miglioramento della gestione di queste eventualità:

\paragraph{Errori di selezione.} 
La selezione manuale introduce la possibilità di errore: l'utente sceglie l'agente sbagliato, 
ottiene una risposta fuori dominio o un rifiuto, e deve ricominciare. 


\paragraph{Richieste ibride.} Nella conversazione naturale, l'utente può passare fluidamente da un argomento 
a un altro: ``Quanti ticket sono aperti oggi? E qual è l'NPS del negozio di Bologna?''. 
Con agenti isolati, questa transizione richiede il cambio manuale di agente. 

La visione di Orion è quella di un \emph{orchestratore intelligente} che elimina questi due limiti. 
L'utente interagisce con un'unica interfaccia che analizza la richiesta, ne determina l'intento e la complessità, 
e la inoltra in modo trasparente all'agente specializzato più appropriato. 
Quando Orion non è sufficientemente sicuro della classificazione, pone una domanda di chiarimento prima di procedere, 
in un meccanismo di disambiguazione \emph{human-in-the-loop} ispirato ai principi dell'\emph{active learning}~\cite{settles2012active,diao2023activeprompting}.

Il pattern architetturale di Orion è quello dell'\emph{orchestratore gerarchico} descritto nella letteratura sui sistemi multi-agente basati su 
LLM~\cite{guo2024multiagent}: un \emph{supervisor agent} coordina $K$ agenti specializzati, 
ciascuno con il proprio dominio di competenza, senza che gli agenti interagiscano direttamente tra loro. 
Rispetto ai pattern cooperativi puri~\cite{wu2023autogen,wang2024moa}, l'architettura gerarchica ha il vantaggio 
della semplicità operativa: il flusso è sempre \emph{utente} $\rightarrow$ \emph{router} $\rightarrow$ \emph{un solo agente} $\rightarrow$ \emph{risposta}, 
senza negoziazioni tra agenti che complicano il debugging e aumentano la latenza.

% =====================================================================
\section{Panoramica architetturale}
\label{sec:orion-panoramica}
% =====================================================================

La Fig.~\ref{fig:orion_architettura} riassume la vista d'insieme.
Una richiesta dell'utente segue un percorso lineare: il \emph{frontend} di Orion 
la inoltra al \emph{backend}, che invoca il \emph{routing step} LLM-based (Sez.~\ref{sec:orion-routing}); 
questo seleziona in un solo passo l'agente competente all'interno dell'insieme disponibile (\texttt{self}, \texttt{sysaid}, \texttt{allycare}, \texttt{coopolicy}). 
La richiesta è quindi inoltrata al backend del sub-agente selezionato tramite l'\emph{agent proxy} (Sez.~\ref{sec:orion-proxy}), 
che ne normalizza la risposta prima di restituirla al frontend.

Il valore \texttt{self} non corrisponde a un sub-agente esterno, 
ma alla pipeline generalista interna di Orion (clonata da Vera, vista in precedenza: RAG sui documenti caricati dall'utente, ricerca web e conoscenza parametrica)
che funge anche da \emph{fallback} quando un sub-agente non è raggiungibile. 
Le sezioni successive approfondiscono i due componenti chiave, il routing step e l'agent proxy, 
preceduti dal posizionamento rispetto ai sistemi affini (Sez.~\ref{sec:orion-related}).

\begin{figure}[htbp]
    \centering
    \includegraphics[width=\textwidth]{orion1}
    \caption[Architettura di Orion]{Architettura dell'orchestratore Orion: il routing step LLM-based seleziona dinamicamente quale agente specializzato invocare; 
    il proxy normalizza le risposte; la pipeline Vera (RAG + web search) resta come fallback su \texttt{agent == "self"}.}
    \label{fig:orion_architettura}
\end{figure}

\begin{figure}[htbp]
    \centering
    \includegraphics[width=0.9\textwidth,height=0.9\textheight,keepaspectratio]{diagrammaconorion}
    \caption[Architettura di sistema con Orion]{La suite nella sua forma finale, all'interno della rete privata di Coop Alleanza: 
    un load balancer interno separa il percorso utente (Orion Frontend $\to$ Orion Backend $\to$ sub-agenti specializzati) dal percorso amministrativo, 
    e tutti i backend condividono l'infrastruttura GCP (Cloud SQL, BigQuery, Secret Manager, Vertex AI con Gemini, Claude, Model Armor e RAG Engine). 
    Evoluzione della Fig.~\ref{fig:architettura_suite} con l'aggiunta dello strato di orchestrazione.}
    \label{fig:architettura_orion_v2}
\end{figure}

% =====================================================================
\section{Lavori correlati e posizionamento}
\label{sec:orion-related}
% =====================================================================

Prima di descrivere l'architettura nel dettaglio, è utile posizionare Orion rispetto ai principali sistemi di routing 
e orchestrazione di LLM proposti nella letteratura recente. 

La Tab.~\ref{tab:orion-related} riassume il confronto lungo cinque dimensioni.

\begin{table}[htbp]
  \centering
  \caption[Confronto con i lavori correlati]{Posizionamento di Orion rispetto ai principali sistemi di routing e orchestrazione multi-agente.}
  \label{tab:orion-related}
  \small
  \begin{tabular}{p{2.4cm}p{3cm}p{2.8cm}p{2.3cm}p{2.3cm}}
    \toprule
    \textbf{Sistema}   & \textbf{Tipo routing} & \textbf{Aggiunta\newline agenti} & \textbf{Confidence} & \textbf{Grounding} \\
    \midrule
    FrugalGPT~\cite{chen2023frugalgpt}
        & Cascade sequenziale    & Riordinamento coda  & Early-exit implicita  & Non previsto \\
    RouteLLM~\cite{ong2024routellm}
        & Classificatore appreso & Riaddestramento     & Non esplicita         & Non previsto \\
    AutoGen~\cite{wu2023autogen}
        & Cooperativo (chat)     & Registrazione agente & Non prevista         & A carico agente \\
    MoA~\cite{wang2024moa}
        & Competitivo + aggr.    & Aggiunta al pool    & Aggregazione          & Non previsto \\
    \textbf{Orion}
        & LLM-based strutturato  & Aggiornamento prompt & Esplicita + clarif.  & Eterogeneo \\
    \bottomrule
  \end{tabular}
\end{table}

\paragraph{Orion vs FrugalGPT.} FrugalGPT~\cite{chen2023frugalgpt} 
formalizza il routing come un \emph{cascade}: i modelli sono ordinati per costo crescente 
e la richiesta scala al modello successivo solo se il precedente non è abbastanza sicuro. 
Orion non è un cascade ma un \emph{router diretto}: il routing step seleziona un singolo agente in un solo passo, senza prove sequenziali. 
Il vantaggio è una latenza prevedibile e indipendente dal numero di agenti; lo svantaggio è che la decisione è irreversibile nel turno.

\paragraph{Orion vs RouteLLM.} RouteLLM~\cite{ong2024routellm} addestra un classificatore su dati 
di preferenza umana per decidere tra un modello forte e uno debole. 
Orion affronta un problema più ampio (selezione tra $K$~agenti eterogenei,
non tra due modelli) e adotta un router LLM-based anziché appreso. Il trade-off è esplicito: 
flessibilità (aggiungere un agente richiede l'aggiornamento del prompt, non il riaddestramento) 
contro efficienza (una chiamata LLM per ogni richiesta). 
La transizione verso un router appreso è discussa come sviluppo futuro (Sez.~\ref{subsec:concl-routellm}).

\paragraph{Orion vs AutoGen.} AutoGen~\cite{wu2023autogen} implementa un pattern cooperativo in cui gli agenti 
dialogano tra loro attraverso un \texttt{GroupChatManager}. Orion adotta un pattern strettamente gerarchico: 
il flusso unidirezionale già descritto (Sez.~\ref{sec:orion-motivazione}) resta l'unica via, 
senza interazione inter-agente. La semplicità del flusso lineare è un requisito in contesto enterprise, 
dove la tracciabilità di ogni decisione e la prevedibilità della latenza sono prioritarie rispetto alla flessibilità compositiva.

\paragraph{Orion vs Mixture-of-Agents.} MoA~\cite{wang2024moa} sottopone la stessa richiesta a $N$~agenti e aggrega le risposte. 
Orion seleziona un singolo agente, evitando il costo computazionale di $N$~invocazioni parallele. 
L'approccio è giustificato dal fatto che gli agenti della suite operano su fonti dati disgiunte: 
aggregare una risposta RAG con una risposta Text-to-SQL non produrrebbe una risposta migliore, ma solo rumore.

\paragraph{Specificità di Orion.} Due caratteristiche distinguono Orion dai sistemi citati: (i)~l'output 
strutturato del router include un \emph{confidence score} esplicito con soglia configurabile e meccanismo di \emph{clarification}, 
implementando una forma di \emph{deferral} basato sull'incertezza, ispirato ai principi dell'\emph{active learning} 
ma operante in fase di inferenza (\emph{human-in-the-loop}) anziché come ciclo di ri-addestramento; (ii)~gli agenti orchestrati utilizzano forme di 
grounding eterogenee (Web, RAG, Text-to-SQL), una dimensione assente nei lavori citati, 
che si concentrano sulla selezione del modello piuttosto che sulla selezione della \emph{fonte di conoscenza}.

% =====================================================================
\section{Il routing step: stima di intento, complessità, confidenza}
\label{sec:orion-routing}
% =====================================================================

Il cuore di Orion è il \emph{routing step}: un modulo che, 
dato il prompt dell'utente e l'eventuale cronologia della conversazione, 
produce una decisione di instradamento sotto forma di oggetto strutturato. 
La scelta di delegare il routing a un LLM, anziché a un classificatore addestrato ad hoc, 
è motivata da tre ragioni pratiche: (i)~il set di agenti target evolve nel tempo 
 e non si desidera riaddestrare un modello a ogni aggiunta; (ii)~il reasoning esplicito del modello 
 fornisce tracciabilità diagnostica per ogni decisione; (iii)~il costo di una singola chiamata 
 a un modello veloce ed economico (Gemini~2.5 Flash) è trascurabile rispetto al costo 
 della risposta completa dell'agente invocato.

\subsection{Output strutturato del router}
\label{subsec:orion-output-strutturato}

Il router produce un oggetto JSON vincolato dallo schema seguente, 
utilizzando il meccanismo di \emph{guided generation}~\cite{willard2023outlines} offerto dall'API di 
Vertex AI (\texttt{response\_mime\_type = "application/json"} con \texttt{response\_schema} esplicito):

\begin{lstlisting}[language=Python,caption={Schema dell'output strutturato del router.},label={lst:router-schema}]
{
  "agent":      "self | sysaid | allycare | coopolicy",
  "complexity": "low | high",
  "confidence": 0.0 ... 1.0,
  "reasoning":  "Chain-of-Thought esplicita",
  "clarification_needed": null | "<domanda all'utente>"
}
\end{lstlisting}

Ogni campo assolve a una funzione precisa:

\begin{description}
  \item[\texttt{agent}] Identificatore dell'agente target. 
  Il campo è vincolato a un'\texttt{enum} contenente i soli agenti disponibili per l'utente corrente, 
  garantendo che il modello non possa inventare target inesistenti. 
  Il valore \texttt{"self"} indica che Orion gestisce la richiesta internamente 
  (RAG su documenti caricati dall'utente, ricerca web, conoscenza parametrica).

  \item[\texttt{complexity}] Stima binaria della complessità della richiesta: \texttt{"low"} per domande dirette, 
  conteggi semplici, filtri base; \texttt{"high"} per analisi multi-dimensionali, correlazioni, 
  confronti temporali articolati, report executive. 
  La complessità non influenza la scelta dell'agente ma è propagata al sub-agente nel 
  payload della richiesta, che la utilizza per il proprio \emph{model routing} interno (cfr. Sez.~\ref{subsec:agenti-sysaid-confronto}), 
  realizzando il trade-off costo--qualità descritto nella Sez.~\ref{sec:bg-routing}.

  \item[\texttt{confidence}] Valore numerico nell'intervallo $[0, 1]$ 
  che esprime la certezza del router nella classificazione effettuata. 
  Il campo è il fondamento del meccanismo di \emph{clarification} descritto nella sezione successiva.

  \item[\texttt{reasoning}] Catena di ragionamento esplicita (\emph{Chain-of-Thought}~\cite{wei2022cot}) 
  che il modello è istruito a produrre prima della decisione. 
  Il campo non viene mostrato all'utente ma è persistito nel log di routing e 
  consultabile dall'amministratore attraverso il pannello di amministrazione, 
  abilitando un ciclo di miglioramento continuo del prompt.

  \item[\texttt{clarification\_needed}] Domanda di chiarimento da porre all'utente quando la confidence è insufficiente. 
  Quando valorizzato, il routing non procede: Orion restituisce la domanda all'utente e 
  attende un nuovo turno conversazionale prima di classificare nuovamente.

\end{description}

\subsection{Logica di confidence e clarification}
\label{subsec:orion-confidence}

La \emph{confidence} prodotta dal router governa il comportamento del 
sistema secondo una logica a soglia configurabile. 
La soglia di default, \texttt{ROUTER\_CONFIDENCE\_THRESHOLD}, è impostata a 0,60 (modificabile):

\begin{itemize}
  \item Se $\text{confidence} \geq \text{threshold}$, il routing procede normalmente: la richiesta è inoltrata all'agente selezionato.
  \item Se $\text{confidence} < \text{threshold}$, il routing è sospeso. Se il modello ha già proposto una domanda di chiarimento (\texttt{clarification\_needed} non nullo), essa viene restituita all'utente. In caso contrario, il sistema genera automaticamente una richiesta di disambiguazione generica.
\end{itemize}

La calibrazione del confidence score, la proprietà per cui il modello è 
più incerto quando sbaglia e più sicuro quando ha ragione, non è banale;
la letteratura mostra che i modelli di grande scala sono ``ragionevolmente'' calibrati quando interrogati
esplicitamente sulla propria certezza~\cite{kadavath2022knowsknows}, ma la calibrazione peggiora su domini
specifici non visti in addestramento. Il Capitolo~\ref{chap:orion-valutazione} verifica empiricamente questa proprietà sul dominio del progetto.

\begin{figure}[htbp]
    \centering
    \includegraphics[height=0.85\textheight,keepaspectratio]{orion2}
    \caption[Routing decision flow]{Logica di routing: in funzione della soglia di confidenza configurabile, Orion instrada direttamente oppure pone all'utente una domanda di chiarimento.}
    \label{fig:routing_decision}
\end{figure}

\subsection{Scelta del modello per il routing}
\label{subsec:orion-model-routing}

Il modello utilizzato per il routing step è Gemini~2.5 Flash~\cite{team2024gemini}. 
La scelta è motivata da tre fattori:

\begin{enumerate}
  \item \textbf{Latenza.} Il routing step aggiunge latenza all'intera pipeline: 
  il tempo trascorso tra la richiesta dell'utente e l'inizio della risposta è la somma della latenza del router e della latenza dell'agente. 
  Con Gemini Flash, il modello più rapido della famiglia, la latenza aggiunta dal routing step è contenuta e trascurabile rispetto al tempo di generazione della risposta completa dell'agente.
  \item \textbf{Costo.} Il router è chiamato su \emph{ogni} richiesta utente. Un modello più potente (ad esempio Gemini~2.5 Pro o Claude Sonnet) avrebbe un costo unitario superiore, 
  che moltiplicato per il volume di traffico rischia di diventare impattante.
  \item \textbf{Capacità sufficiente.} Il task del router è una classificazione 
  vincolata su un set chiuso di $K$~etichette (quattro agenti + la variabile di complessità), con ragionamento esplicito. 
  Un modello ``flash'' dovrebbe sempre essere sufficiente per questo tipo di compito,
  come confermato dai risultati sperimentali del Capitolo~\ref{chap:orion-valutazione}.
\end{enumerate}

L'alternativa di un router appreso, come RouteLLM~\cite{ong2024routellm} (Sez.~\ref{sec:bg-routing}), 
è stata valutata come possibile evoluzione futura una volta che il sistema 
sarà passato in produzione a tutti gli effetti, con l'aumento del volume di traffico
che potrebbe generare un dataset di training significativo.

% =====================================================================
\section{Agent proxy e normalizzazione delle risposte}
\label{sec:orion-proxy}
% =====================================================================

Una volta che il router ha selezionato l'agente target, il \emph{proxy} di Orion inoltra 
la richiesta al backend del sub-agente via HTTP e normalizza la risposta prima di restituirla al frontend. 
Il proxy è il componente che rende trasparente al frontend l'eterogeneità dei sub-agenti: 
ogni agente ha il proprio formato di risposta (SysAid include query SQL e grafici base64; CooPolicy include citazioni con URI ai documenti; AllyCare include tabelle dati), 
e il proxy li mappa verso una \emph{shape} uniforme.

Le decisioni architetturali rilevanti del proxy sono quattro.

\paragraph{Nessuna modifica ai backend dei sub-agenti.} Il proxy si adatta al contratto API esistente di ogni sub-agente, 
senza richiedere modifiche ai loro endpoint. Questo principio preserva la stabilità degli agenti già in produzione e la compatibilità.

\paragraph{Autenticazione service-to-service.} I sub-agenti non ricevono il token Entra ID dell'utente: 
L'autenticazione fra i backend è gestita semplicemente fornendo al service account del backend di Orion il permesso di invocare gli altri servizi.
Le informazioni dell'utente (nome, ruolo, direzione), necessarie ai sub-agenti per personalizzare le risposte, 
sono propagate come campi nel payload JSON della richiesta. 
Questo approccio elimina la eventuale complessità del \emph{token forwarding}.

\paragraph{Normalizzazione delle risposte.} Il proxy riceve dalla risposta del sub-agente un JSON 
con campi eterogenei e lo trasforma in un oggetto \texttt{OrionProxyResponse} con una struttura fissa. 
Un caso particolarmente significativo è la normalizzazione delle citazioni di CooPolicy: gli URI dei documenti, 
originariamente nel formato \texttt{gs://bucket/file.pdf}, vengono riscritti come \texttt{/api/orion/docs/file?uri=...}, 
un percorso relativo che il frontend può aprire direttamente e che viene servito da un proxy autenticato verso GCS montato nel backend di Orion.

\paragraph{Gestione degli errori e fallback.} Se il sub-agente risponde con un codice di errore (4xx, 5xx) 
o non risponde entro il timeout di 120~secondi, il proxy restituisce un messaggio di errore leggibile e il flag \texttt{error = true}. 
Il frontend rileva il flag e può mostrare all'utente un messaggio appropriato. In alternativa, 
il flusso può ricadere sulla pipeline interna di Orion (\texttt{agent = "self"}).

\paragraph{Streaming bufferizzato.} Il proxy attende la risposta completa del sub-agente 
prima di inoltrarla al frontend via \emph{Server-Sent Events} (SSE). 
La scelta semplifica la normalizzazione al costo di un \emph{time-to-first-token}
pari alla latenza dell'intero sub-agente; l'evoluzione verso uno streaming \emph{pass-through} è una possibile evolutiva architetturale (Tab.~\ref{tab:orion-decisioni}).

\begin{figure}[htbp]
    \centering
    \includegraphics[width=\textwidth]{orion3}
    \caption[Sequence diagram di un turno Orion]{Sequence diagram completo: dalla richiesta utente al routing, dall'invocazione del sub-agente alla normalizzazione della risposta.}
    \label{fig:orion_sequence}
\end{figure}

% =====================================================================
\section{Conversation threading e gestione della storia}
\label{sec:orion-history}
% =====================================================================

Orion mantiene una \emph{storia unica} per conversazione: tutti i turni, indipendentemente dall'agente che li ha serviti, sono memorizzati nella stessa conversazione del database di Orion. 
Questa scelta consente all'utente di cambiare argomento (e quindi agente) all'interno 
della stessa conversazione senza discontinuità percepita.

Il trade-off accettato è che il sub-agente invocato riceve l'intera storia della conversazione, 
inclusi i turni gestiti da altri agenti. Per mitigare il rischio di confusione, la storia inviata al sub-agente 
è troncata agli ultimi sei turni (parametro configurabile) e i messaggi non pertinenti (ad esempio, grafici base64 o citazioni di documenti) 
sono sostituiti con un riassunto testuale sintetico. Il router stesso riceve la storia per 
migliorare la classificazione dei turni di follow-up: se l'utente ha appena chiesto informazioni sui ticket IT, 
una domanda successiva come ``E quanti ne sono stati chiusi?'' viene correttamente instradata a SysAid anche in assenza di parole chiave esplicite.

% =====================================================================
\section{Decisioni architetturali e loro motivazione}
\label{sec:orion-decisioni}
% =====================================================================

Questa sezione riassume le principali decisioni architetturali di Orion, 
documentando per ciascuna le alternative considerate e le ragioni della scelta effettuata.

\begin{table}[htbp]
  \centering
  \caption[Decisioni architetturali di Orion]{Principali decisioni architetturali di Orion con alternative considerate e motivazione della scelta.}
  \label{tab:orion-decisioni}
  \begin{tabular}{p{3.5cm}p{3.5cm}p{6cm}}
    \toprule
    \textbf{Decisione}       & \textbf{Alternativa}        & \textbf{Motivazione}                   \\
    \midrule
    Modello routing: Gemini Flash
                             & Gemini Pro, Claude, router appreso~\cite{ong2024routellm}
                             & Latenza sub-secondo, costo trascurabile, capacità sufficiente per classificazione a $K$ etichette \\
    Soglia confidence: 0,60  & 0,50, 0,70, adattiva
                             & Compromesso tra clarification rate e accuracy, calibrato empiricamente (Cap.~\ref{chap:orion-valutazione}) \\
    Storia: unica in Orion   & Separata per agente
                             & Continuità dell'esperienza utente; i sub-agenti ricevono storia troncata \\
    Routing trasparente      & Nascosto
                             & Badge visibile sopra la risposta con agente e complessità; valore dimostrativo e fiducia utente \\
    \bottomrule
  \end{tabular}
\end{table}
